\documentclass[aspectratio=169]{beamer}

\usetheme{Madrid}
\usecolortheme{default}
\setbeamertemplate{navigation symbols}{}
\setbeamertemplate{footline}[frame number]
\setbeamertemplate{section page}{
  \centering
  \vfill
  {\usebeamerfont{section title}\insertsection\par}
  \vfill
}

\usepackage[utf8]{inputenc}
\usepackage[T1]{fontenc}
\usepackage{amsmath, amssymb, bm}
\usepackage{booktabs}
\usepackage{hyperref}

\title[Research Meeting]{September 11, 2026 Research Meeting}
\subtitle{}
\author{Siyan Wang}
\date{}

\begin{document}

\begin{frame}
  \titlepage
\end{frame}

\begin{frame}{What I did this week}
  \begin{itemize}
    \item Understood fingerprint-based detection and attribution methods.
    \item Read and summarized Naveau et al. (2020) on extreme event attribution.
    \item Found the topic (2026 European heat wave) for case study and its relative reports.
    \item Prepared slides for the research meeting.
  \end{itemize}
  
\end{frame}


\section{Fingerprinting}

\begin{frame}[plain]
  \sectionpage
\end{frame}

\begin{frame}{What Is Fingerprinting?}
\small
\begin{itemize}
  \item In general, a fingerprint is a distinctive pattern used to identify or attribute a source.
  \item In climate science, a fingerprint is the expected response pattern to a specific forcing.
  \item Examples of forcings:
  \[
    GHG,\quad aerosols,\quad natural\ forcing,\quad ALL.
  \]
  \item The attribution question is whether the observed climate record contains this expected pattern.
  \item Fingerprinting is usually used for long-term detection and attribution.
\end{itemize}
\end{frame}

\begin{frame}{Fingerprinting as a Signal-plus-Noise Problem}
\small
A simplified detection and attribution model is:
\[
  \mathbf{y}
  =
  \sum_{k=1}^{K} \beta_k \mathbf{x}_k
  +
  \boldsymbol{\epsilon}.
\]

\begin{itemize}
  \item $\mathbf{y}$: observed climate change pattern.
  \item $\mathbf{x}_k$: model-simulated fingerprint for forcing $k$.
  \item $\beta_k$: scaling factor.
  \item $\boldsymbol{\epsilon}$: internal climate variability and residual error.
\end{itemize}

\[
  \beta_k > 0
  \quad \Rightarrow \quad
  \text{signal detected.}
\]
\end{frame}

\begin{frame}{Optimal Fingerprinting}
\small
Optimal fingerprinting estimates the scaling factors while accounting for the covariance of internal variability.

\[
  \mathbf{y}
  =
  \mathbf{X}\boldsymbol{\beta}
  +
  \boldsymbol{\epsilon},
  \qquad
  \boldsymbol{\epsilon}\sim N(\mathbf{0},\mathbf{C}).
\]

The generalized least squares estimator is:
\[
  \hat{\boldsymbol{\beta}}
  =
  (\mathbf{X}^{T}\mathbf{C}^{-1}\mathbf{X})^{-1}
  \mathbf{X}^{T}\mathbf{C}^{-1}\mathbf{y}.
\]

\begin{itemize}
  \item $\mathbf{C}$ is estimated from control runs or internal-variability simulations.
  \item $\mathbf{C}^{-1}$ weights directions with lower noise more strongly.
  \item This improves signal-to-noise separation.
\end{itemize}
\end{frame}

\begin{frame}{Use in Nested-Model Comparison}
\small
Fingerprinting can be framed as comparing nested regression models.

\[
  H_0: \mathbf{y} = \mathbf{X}_{-k}\boldsymbol{\beta}_{-k} + \boldsymbol{\epsilon}
\]

\[
  H_1: \mathbf{y} = \mathbf{X}_{-k}\boldsymbol{\beta}_{-k}
  + \beta_k \mathbf{x}_k + \boldsymbol{\epsilon}
\]

\begin{itemize}
  \item $H_0$: the forcing-$k$ fingerprint is excluded.
  \item $H_1$: the forcing-$k$ fingerprint is included.
  \item If adding $\mathbf{x}_k$ significantly improves the fit, the signal is detected.
  \item This is usually tested through confidence intervals for $\beta_k$ or likelihood / residual-sum-of-squares comparisons.
\end{itemize}

\[
  \beta_k = 0
  \quad \text{vs.} \quad
  \beta_k \ne 0
\]
\end{frame}

\section{Statistical Methods for Extreme Event Attribution in Climate Science}

\begin{frame}[plain]
  \sectionpage
  \vspace{-1.2cm}
  \begin{center}
    {\large Naveau, Hannart, and Ribes (2020)\par}
    \vspace{0.25cm}
    {\small Annual Review of Statistics and Its Application, 7:89--110\par}
  \end{center}
\end{frame}

\begin{frame}{Why This Paper Matters}
  \begin{itemize}
    \item Extreme event attribution (EEA) asks whether human influence changed the likelihood or intensity of a class of extreme events.
    \item The paper is a topic review: it maps the statistical workflow, assumptions, and open problems.
    \item Its key message: EEA is not just a climate-model exercise; it is a probability, uncertainty, tail-risk, and causal-inference problem.
  \end{itemize}
\end{frame}

\begin{frame}{The EEA Question}
  \begin{block}{Common public question}
    Did climate change cause this heat wave, flood, drought, or record event?
  \end{block}

  \begin{block}{Statistical EEA question}
    Did anthropogenic forcing change the probability of this \emph{class of events}?
  \end{block}

  \vspace{0.2cm}
  \begin{itemize}
    \item EEA therefore defines a repeatable event class: variable, region, time window, and threshold.
    \item The final attribution statement is conditional on this definition.
  \end{itemize}
\end{frame}

\section{Section 1: Introduction}

\begin{frame}{Section 1: Starting Point}
  \begin{itemize}
    \item EEA compares extreme-event likelihoods under different scenarios.
    \item The two central worlds are:
      \begin{itemize}
        \item factual world: close to the observed climate, including anthropogenic forcing;
        \item counterfactual world: a hypothetical climate without anthropogenic emissions.
      \end{itemize}
    \item The paper focuses on statistical methods rather than the storyline approach.
  \end{itemize}
\end{frame}

\begin{frame}{Workflow Lesson from Section 1}
  \begin{enumerate}
    \item Start from a real impact or observed extreme.
    \item Translate it into a statistical event class.
    \item Define factual and counterfactual climates.
    \item Estimate and compare event probabilities.
  \end{enumerate}

  \vspace{0.25cm}
  \begin{alertblock}{Practical warning}
    The event definition is not a technical afterthought. It determines what question the attribution result actually answers.
  \end{alertblock}
\end{frame}

\section{Section 2: Basic Setups}

\begin{frame}{Section 2: Probability Setup}
  Let $u$ be a high threshold and let the event be an exceedance.

  \begin{align*}
    p_1(u) &= P(Z > u) && \text{factual world}, \\
    p_0(u) &= P(X > u) && \text{counterfactual world}.
  \end{align*}

  \begin{itemize}
    \item $Z$ and $X$ represent the same climate quantity in two different worlds.
    \item Example: regional  3-day Annual/June maximum temperature mean (Tx3x) in 
    current climate verse a simulation without anthropogenic or previous year's forcing (GMST).
  \end{itemize}
\end{frame}

\begin{frame}{Risk Ratio and FAR}
  Two common attribution quantities are:

  \begin{align*}
    RR(u) &= \frac{p_1(u)}{p_0(u)}, \\
    FAR(u) &= 1 - \frac{p_0(u)}{p_1(u)}.
  \end{align*}

  \begin{itemize}
    \item $RR=2$: the event is twice as likely in the factual world.
    \item $FAR=0.5$: half of the event probability is attributable to the forcing, under the assumed framework.
    \item Both depend on the threshold $u$.
  \end{itemize}
\end{frame}

\begin{frame}{How RR Is Computed in Practice}
  \begin{enumerate}
    \item Define $X$ and $u$: for example, June regional Tx3x exceeding the observed event value.
    \item Obtain samples from factual and counterfactual worlds.
    \item Estimate exceedance probabilities:
      \[
        \hat p_1 = \frac{\#\{Z_i > u\}}{n_1}, \qquad
        \hat p_0 = \frac{\#\{X_i > u\}}{n_0}.
      \]
    \item Compute:
      \[
        \widehat{RR} = \frac{\hat p_1}{\hat p_0}.
      \]
  \end{enumerate}
\end{frame}

\begin{frame}{Why Section 2 Is Not Just a Formula}
  \begin{itemize}
    \item If $p_0$ is very small, simple counting can give zero exceedances.
    \item Then $RR$ becomes unstable or effectively unbounded.
    \item Return periods are harder in a nonstationary climate: a ``100-year event'' depends on the climate state.
    \item This motivates later sections on data design, event definition, and extreme value theory.
  \end{itemize}
\end{frame}

\begin{frame}{Figure 2: Causality Depends on Threshold}

\begin{columns}[T,onlytextwidth]

  % Left: text
  \begin{column}{0.57\textwidth}
    \small
    \begin{align*}
      FAR(u) &= 1-\frac{p_0(u)}{p_1(u)}, \\
      PS(u) &= \max\left\{1-\frac{1-p_1(u)}{1-p_0(u)},0\right\}, \\
      PNS(u) &= \max\{p_1(u)-p_0(u),0\}.
    \end{align*}

    \begin{itemize}
      
      \item As $u$ increases, estimating $p_0$ and $RR$ becomes harder
      because counterfactual exceedances become very rare.

      \item FAR/ $PN$  often increases: if the event occurs, human
      influence is more likely to have been necessary.

      \item $PS$ often decreases: human influence alone is less likely
      to be sufficient for a very extreme event.

      \item $PNS$ is often largest at an intermediate threshold, where
      $p_1(u)-p_0(u)$ is largest.
    \end{itemize}
  \end{column}

  % Right: image
  \begin{column}{0.40\textwidth}
    \centering
    \includegraphics[
      width=\linewidth,
      height=0.72\textheight,
      keepaspectratio
    ]{image/naveau_fig2.png}

    \vspace{0.1cm}
    {\scriptsize Figure 2: Causal probabilities as functions of threshold $u$.}
    
    {\scriptsize Figure 2 shows how these causal probabilities change as
      the threshold $u$ changes.}
  \end{column}

\end{columns}

\end{frame}

\section{Section 3: Data and Uncertainty}

\begin{frame}{Section 3: Sources of Uncertainty}
  Naveau et al. emphasize that factual, counterfactual, and observed climate data are all uncertain.

  \begin{itemize}
    \item natural internal variability;
    \item numerical model uncertainty;
    \item observational uncertainty;
    \item sampling uncertainty in space and time;
    \item statistical model uncertainty.
  \end{itemize}
\end{frame}

\begin{frame}{Data Types in EEA}
  \begin{tabular}{p{0.35\linewidth}p{0.55\linewidth}}
    \toprule
    Data source & Main role and limitation \\
    \midrule
    Conditional simulations & Large samples for one event, but depend on boundary-condition choices. \\
    CMIP transient simulations & Multi-model and broadly available, but smaller ensembles and coarser resolution. \\
    Observations/reanalysis & Anchors the event in reality, but records are short for extremes. \\
    \bottomrule
  \end{tabular}
\end{frame}

\begin{frame}{Workflow Lesson from Section 3}
  Before estimating $RR$, ask:
  \begin{itemize}
    \item Which data represent the factual world?
    \item Which data represent the counterfactual world?
    \item Is the ensemble large enough for the event rarity?
    \item Does the model reproduce the relevant variable, region, and dependence structure?
    \item Which uncertainties are propagated into the final interval?
  \end{itemize}
\end{frame}

\section{Section 4: Event Definitions}

\begin{frame}{Section 4: Defining the Event}
  An event definition maps a spatio-temporal field into a probability statement.

  \[
    Y(t,s) \quad \longrightarrow \quad X = g\{Y(t,s)\} \quad \longrightarrow \quad \{X > u\}.
  \]

  \begin{itemize}
    \item $Y(t,s)$: original climate field.
    \item $X$: summary of the event, such as regional mean heat or maximum rainfall.
    \item $u$: threshold, often the observed event magnitude or a return-level threshold.
  \end{itemize}
\end{frame}

\begin{frame}{The Event Definition Trade-Off}
  \begin{itemize}
    \item Too narrow: physically close to the impact, but low statistical power.
    \item Too broad: easier to estimate, but may lose the event people care about.
    \item Record events are different from fixed-threshold exceedances.
    \item Compound events require multivariate definitions.
  \end{itemize}

  \vspace{0.2cm}
  \begin{alertblock}{Practical warning}
    The event definition should be physically meaningful and statistically estimable.
  \end{alertblock}
\end{frame}

\begin{frame}{Advanced Event Definition Ideas}
  \small
  The paper discusses methods that make event definition part of the statistical problem.

  \begin{block}{Cattiaux and Ribes: most exceptional domain}
    Instead of fixing the region and duration in advance, search over candidate
    regions $R$ and durations $D$ to find where the observed event is most
    exceptional. \\
    
    Example: for a heat wave, compare 3-day, 10-day, monthly, or
    seasonal definitions over different spatial domains.
  \end{block}

  \begin{block}{Hannart and Naveau: most causal projection}
    Treat the event as a multivariate field and find a projection, that carries strong counterfactual causal information, for
    example by optimizing a quantity related to $PNS$.
  \end{block}

  \vspace{-0.05cm}
  \begin{itemize}
    \item Benefit: the event definition is more targeted to attribution.
    \item Cost: one must choose the search space, optimization criterion,
    projection form, and how to correct for selection bias.
  \end{itemize}
\end{frame}

\section{Section 5: Statistical Inference}

\begin{frame}{Section 5: The Statistical Task}
  Once the event is defined, the core task is to estimate two rare-event probabilities:
  \[
    p_1(u)=P(Z>u), \qquad p_0(u)=P(X>u).
  \]

  \vspace{0.15cm}
  \begin{itemize}
    \item $Z$: the event variable in the factual world.
    \item $X$: the same event variable in the counterfactual world.
    \item $u$: the event threshold, often the observed event magnitude.
  \end{itemize}

  \vspace{0.15cm}
  The paper reviews four directions: counting, EVT, hierarchical modeling, and multivariate extremes.
\end{frame}

\begin{frame}{Nonparametric Counting: The Simple Estimator}
  Suppose we have factual samples $Z_1,\ldots,Z_n$ and counterfactual samples
  $X_1,\ldots,X_n$.

  \[
    \hat p_1(u)=\frac{1}{n}\sum_{i=1}^n I(Z_i>u),
    \qquad
    \hat p_0(u)=\frac{1}{n}\sum_{i=1}^n I(X_i>u).
  \]

  \[
    \widehat{RR}(u)=\frac{\hat p_1(u)}{\hat p_0(u)}, \qquad
    \widehat{FAR}(u)=1-\frac{\hat p_0(u)}{\hat p_1(u)}.
  \]

  \begin{itemize}
    \item Strength: transparent and distribution-free.
    \item Requirement: the ensemble must be much larger than the return period of interest.
  \end{itemize}
\end{frame}

\begin{frame}{Why Counting Can Fail}
  Naveau et al. show that the uncertainty of the counting estimator grows quickly in the tail.

  \[
    \sqrt{n}\{\widehat{FAR}(u)-FAR(u)\}
    \approx N\{0,\sigma^2(u)\}.
  \]

  \[
    \sigma^2(u)
    =
    \frac{p_0^2(u)}{p_1^2(u)}
    \left[
      \frac{1-p_0(u)}{p_0(u)}
      +
      \frac{1-p_1(u)}{p_1(u)}
    \right].
  \]

  \begin{itemize}
    \item If $p_0(u)$ is tiny, counterfactual exceedances may be zero.
    \item Then $RR$ becomes unstable, even if the causal signal is strong.
  \end{itemize}
\end{frame}

\begin{frame}{EVT: Modeling the Tail}
  EVT replaces raw counting with a tail model.

  \begin{block}{Block maxima}
    If $X$ is a block maximum, such as annual maximum daily temperature,
    use the generalized extreme value distribution:
    \[
      X \sim GEV(\mu,\sigma,\xi).
    \]
  \end{block}

  \begin{block}{Threshold exceedances}
    If we model exceedances above a high cutoff $v$, use a generalized Pareto tail:
    \[
      P(X>x \mid X>v) \approx H_{\xi}\left(\frac{x-v}{\sigma}\right),
      \qquad x>v.
    \]
  \end{block}

  Parameters are commonly estimated by maximum likelihood.
\end{frame}

\begin{frame}{Computing RR with a GEV Model}
  For a heat wave index, define a yearly extreme:
  \[
    Y_t = \text{annual or seasonal maximum 3-day mean temperature}.
  \]

  A common EEA model is nonstationary:
  \[
    Y_t \sim GEV(\mu_t,\sigma,\xi),
    \qquad
    \mu_t=\mu_0+\alpha \, GMST_t.
  \]

  For event threshold $u$:
  \[
    p_1 = 1-F_{GEV}(u;\mu_{\text{factual}},\sigma,\xi),
  \]
  \[
    p_0 = 1-F_{GEV}(u;\mu_{\text{counterfactual}},\sigma,\xi),
    \qquad
    RR=\frac{p_1}{p_0}.
  \]
\end{frame}

\begin{frame}{From Estimated Parameters to RR}
  \small
  After fitting the GEV, plug in two climate states:

  \[
    \hat\mu_1=\hat\mu_0+\hat\alpha\,GMST_{\text{factual}},
    \qquad
    \hat\mu_{cf}=\hat\mu_0+\hat\alpha\,GMST_{\text{counterfactual}}.
  \]

  Then compute fitted tail probabilities:
  \[
    \hat p_1
    =
    1-F_{GEV}(u;\hat\mu_1,\hat\sigma,\hat\xi),
    \qquad
    \hat p_0
    =
    1-F_{GEV}(u;\hat\mu_{cf},\hat\sigma,\hat\xi).
  \]

  \[
    \widehat{RR}=\frac{\hat p_1}{\hat p_0}.
  \]

  Uncertainty is usually obtained by bootstrap, profile likelihood, Bayesian posterior sampling, or delta-method approximations.
\end{frame}

\begin{frame}{EVT: Benefit and Cost}
  \begin{itemize}
    \item Benefit: EVT can estimate probabilities beyond the observed or simulated sample range.
    \item This is crucial when the counterfactual event probability $p_0$ is extremely small.
    \item Cost: the result depends on tail assumptions: GEV/GPD choice, threshold choice, stationarity, and tail shape $\xi$.
    \item In EEA, model misspecification in the tail can strongly distort $RR$ and FAR.
  \end{itemize}
\end{frame}

\begin{frame}{Hierarchical and Multivariate Extensions}
  The paper argues that richer models are needed when one data set and one tail model are too narrow.

  \begin{itemize}
    \item Hierarchical models can combine observations, climate models, internal variability, and model uncertainty.
    \item Multivariate EVT is needed for compound events, such as heat plus drought or rainfall plus saturated soil.
    \item Tail dependence matters: wrong assumptions about dependence can seriously misestimate joint risk.
    \item These directions are promising, but still underdeveloped in routine EEA.
  \end{itemize}
\end{frame}

\section{Section 6: Conclusions and Open Problems}

\begin{frame}{Section 6: Main Challenges}
  Naveau et al. identify several statistical challenges:

  \begin{itemize}
    \item combine all known sources of uncertainty;
    \item define relevant spatial and temporal event scales;
    \item estimate extremely small probabilities;
    \item handle compound and concomitant extremes;
    \item address model misspecification and spatial multiple testing;
    \item connect statistical attribution with physical causal storylines.
  \end{itemize}
\end{frame}

\begin{frame}{A Practical EEA Workflow}
  \begin{enumerate}
    \item State the attribution question.
    \item Define factual and counterfactual worlds.
    \item Define the event: variable, region, duration, threshold.
    \item Choose data: observations, reanalysis, CMIP, or conditional ensembles.
    \item Identify uncertainty sources.
    \item Estimate $p_1$ and $p_0$.
    \item Compute $RR$, FAR, or causality-related quantities.
    \item Quantify uncertainty.
    \item Interpret as probability change, not deterministic causation.
  \end{enumerate}
\end{frame}

\begin{frame}{Key Takeaways}
  \begin{itemize}
    \item EEA is a comparison of factual and counterfactual event probabilities.
    \item The hard part is not the ratio itself, but credible estimation of $p_1$ and $p_0$.
    \item Event definition controls the scientific meaning of the result.
    \item Rare events require careful tail modeling and uncertainty quantification.
    \item Future EEA needs stronger integration of climate modeling, EVT, hierarchical statistics, and causal inference.
  \end{itemize}
\end{frame}


\begin{frame}{Reference}
  \small
  Naveau, P., Hannart, A., and Ribes, A. (2020). Statistical Methods for Extreme Event Attribution in Climate Science. \emph{Annual Review of Statistics and Its Application}, 7, 89--110.
\end{frame}

\end{document}
